\section{Experiments}
\label{sec:experiments}

All experiments reuse the human-labeled comparison dataset provided for this project. We run the
judge on every comparison and use $n$ human labels, with $n/N\in\{1,2,5,10,20\}\%$. A test set split
by prompt is used only for evaluation.

\paragraph{E1: Estimation.}
We compare \dfsp{} with the following estimators:
\begin{itemize}
  \item the judge's win probability;
  \item a kernel estimator using human labels only;
  \item Platt scaling and isotonic regression of $p_J$;
  \item the linear map of \citet{vandenburg2025aligning};
  \item a GAM and boosted trees on $Z$;
  \item \sfsp.
\end{itemize}
We report Brier score, log-loss and calibration error as functions of $n$. We also run \dfsp{} with
only the tie probability or only the order dependence added to $p_J$.

\paragraph{E2: Does the judge's distribution add information?}
We estimate $G$ from Proposition~\ref{prop:value} and test $G=0$ with the method of
\citet{williamson2023general}. The test is first checked on simulated data with known $G$.

\paragraph{E3: Label selection.}
We compare the rule of Section~\ref{sec:selection} with the following alternatives:
\begin{itemize}
  \item random selection;
  \item selection by judge uncertainty;
  \item Fisher-information selection \citep{shen2025active};
  \item selection of likely judge errors, as in \citet{xu2025rlthf}.
\end{itemize}
We run this comparison on simulated data with known $p_H$ and $\delta$, and on the real data.

\paragraph{E4: Unreliable judge.}
We corrupt the judge's probabilities in four ways: adding a bias, making them overconfident,
reversing them in part of $\cZ$, and adding noise. We check that the selected $\theta_1$ decreases
and that \dfsp{} approaches the human-only estimator.

\paragraph{E5: Post-training.}
We fine-tune a small open-weight model with LoRA \citep{hu2022lora} under the following training
targets:
\begin{itemize}
  \item the $n$ human labels only;
  \item hard judge labels;
  \item judge probabilities, trained with soft DPO and with GAPO \citep{furuta2024gapo};
  \item linear-map estimates;
  \item \dfsp{} estimates;
  \item all human labels, as a reference.
\end{itemize}
We report the held-out human log-loss of $\sigma(d_\vartheta)$, which is the quantity in
Proposition~\ref{prop:bridge}. We also report the win rate against $\pi_{\mathrm{ref}}$, scored by a
reward model trained on all human labels \citep{gao2023scaling} and by a second judge. Finally, we
report the KL divergence to $\pi_{\mathrm{ref}}$ and the response length.

\paragraph{Possible outcomes.}
If $G$ is close to zero, \dfsp{} should perform like \sfsp, and the value of the method lies in the
correction and the label selection alone. If the judge is poor, \dfsp{} should perform like the
human-only estimator. Both outcomes will be reported.
